% !TEX root = ../main.tex
\section{Justificación de cada uno de los campos del registro}

Todos los campos de tipo texto se implementan como \texttt{char[]} de
tamaño fijo, nunca como cadenas dinámicas, porque el tamaño fijo es lo que
permite calcular el \texttt{offset} de cualquier registro dentro de su
archivo mediante \texttt{posición} $\times$ \texttt{sizeof(struct)}.

Cada arreglo de caracteres reserva un byte adicional para el carácter nulo
(\texttt{'\textbackslash 0'}), que en C marca el final real de la cadena:
funciones como \texttt{strlen}, \texttt{strcmp} o
\texttt{printf("\%s", \ldots)} no conocen la longitud de una cadena de
antemano, y sin ese byte seguirían leyendo memoria más allá del dato
válido.

\begin{longtable}{p{2.7cm} p{2.2cm} p{1.5cm} p{6.5cm}}
\toprule
Campo & Tipo & Tamaño & Justificación \\
\midrule
\endhead

carne &
int &
4 B &
Llave primaria numérica, requerida explícitamente por el caso de estudio
junto con la identificación como campo de búsqueda independiente. \\

identificacion &
char[13] &
13 B &
Cubre el caso más ancho entre cédula física (9 dígitos), DIMEX de
residencia (entre 11 y 12 dígitos) y pasaporte alfanumérico: 12 caracteres
más el terminador nulo. \\

tipoIdentificacion &
enum (int) &
4 B &
Distingue nacional/DIMEX/pasaporte sin mezclar esa información dentro del
campo de texto. \\

nombre &
char[21] &
21 B &
20 caracteres útiles, margen razonable para nombres compuestos. \\

apellido1, apellido2 &
char[21] c/u &
42 B &
Separados en vez de un solo campo \texttt{apellidos}, para permitir
búsquedas y ordenamientos por apellido de forma consistente. \\

fechaNacimiento &
struct Fecha &
12 B &
TAD propio (día, mes, año como \texttt{int}), reutilizable y con sus
propias reglas de validez. \\

sexo &
enum (int) &
4 B &
Dominio cerrado, representado como entero en vez de texto libre. \\

carrera &
int &
4 B &
Referencia a un catálogo de carreras en vez de repetir el nombre completo
en cada registro. \\

nivelAcademico &
enum (int) &
4 B &
Dominio cerrado (pregrado, grado, posgrado, etc.). \\

creditosAprobados &
int &
4 B &
Cantidad entera, sin componente fraccionario. \\

promedio &
float &
4 B &
Requiere parte decimal. \\

estado &
int (bool a mano) &
4 B &
0 para inactivo, 1 para activo; sostiene el borrado lógico sin usar
\texttt{bool} de C++. \\

correo &
char[41] &
41 B &
Margen suficiente para direcciones institucionales y personales comunes. \\

telefono &
char[9] &
9 B &
8 dígitos (formato costarricense) más el terminador nulo. \\

direccion &
char[81] &
81 B &
Texto libre de longitud variable, con un tope generoso fijo. \\

\bottomrule
\end{longtable}
